\documentclass[10pt,a4paper]{amsart}
\usepackage{amsmath,amssymb,amsthm,graphicx}
% For commutative diagrams we use Xy-pic
\usepackage[all]{xy}

% We use 2cell for 2-commutative diagrams.
\xyoption{2cell}
% Theorem environments.
%
\theoremstyle{plain}
\newtheorem{theorem}[subsection]{Theorem}
\newtheorem{proposition}[subsection]{Proposition}
\newtheorem{lemma}[subsection]{Lemma}

\theoremstyle{definition}
\newtheorem{definition}[subsection]{Definition}
\newtheorem{example}[subsection]{Example}
\newtheorem{exercise}[subsection]{Exercise}
\newtheorem{situation}[subsection]{Situation}

\theoremstyle{remark}
\newtheorem{remark}[subsection]{Remark}
\newtheorem{remarks}[subsection]{Remarks}

\numberwithin{equation}{subsection}

% Macros
%
\def\lim{\mathop{\mathrm{lim}}\nolimits}
\def\colim{\mathop{\mathrm{colim}}\nolimits}
\def\Spec{\mathop{\mathrm{Spec}}}
\def\Hom{\mathop{\mathrm{Hom}}\nolimits}
\def\Ext{\mathop{\mathrm{Ext}}\nolimits}
\def\SheafHom{\mathop{\mathcal{H}\!\mathit{om}}\nolimits}
\def\SheafExt{\mathop{\mathcal{E}\!\mathit{xt}}\nolimits}
\def\Sch{\mathit{Sch}}
\def\Mor{\mathop{\mathrm{Mor}}\nolimits}
\def\Ob{\mathop{\mathrm{Ob}}\nolimits}
\def\Sh{\mathop{\mathit{Sh}}\nolimits}
\def\NL{\mathop{N\!L}\nolimits}
\def\CH{\mathop{\mathrm{CH}}\nolimits}
\def\proetale{{pro\text{-}\acute{e}tale}}
\def\etale{{\acute{e}tale}}
\def\QCoh{\mathit{QCoh}}
\def\Ker{\mathop{\mathrm{Ker}}}
\def\Im{\mathop{\mathrm{Im}}}
\def\Coker{\mathop{\mathrm{Coker}}}
\def\Coim{\mathop{\mathrm{Coim}}}

% Boxtimes
%
\DeclareMathSymbol{\boxtimes}{\mathbin}{AMSa}{"02}

%
% Macros for moduli stacks/spaces
%
\def\QCohstack{\mathcal{QC}\!\mathit{oh}}
\def\Cohstack{\mathcal{C}\!\mathit{oh}}
\def\Spacesstack{\mathcal{S}\!\mathit{paces}}
\def\Quotfunctor{\mathrm{Quot}}
\def\Hilbfunctor{\mathrm{Hilb}}
\def\Curvesstack{\mathcal{C}\!\mathit{urves}}
\def\Polarizedstack{\mathcal{P}\!\mathit{olarized}}
\def\Complexesstack{\mathcal{C}\!\mathit{omplexes}}
% \Pic is the operator that assigns to X its picard group, usage \Pic(X)
% \Picardstack_{X/B} denotes the Picard stack of X over B
% \Picardfunctor_{X/B} denotes the Picard functor of X over B
\def\Pic{\mathop{\mathrm{Pic}}\nolimits}
\def\Picardstack{\mathcal{P}\!\mathit{ic}}
\def\Picardfunctor{\mathrm{Pic}}
\def\Deformationcategory{\mathcal{D}\!\mathit{ef}}
\usepackage{cleveref}
\crefname{originalreference}{original reference}{original references}
\Crefname{originalreference}{Original reference}{Original references}
\makeatletter
\newlabel{theorem-colimit}{{09YQ}{source}{}{sourcepointer-1}{}}
\newlabel{theorem-colimit@cref}{{[originalreference][0][]09YQ}{[1][0][]source}{}{}{}}
\makeatother
\title{Two formulations of the stalk comparison of cohomology over strict henselisations and a separably closed field}
\author{The Stacks Project Authors}
\date{stacks-project:etale-cohomology.tex:local-snapshot-20261011; GNU FDL 1.2 or later}
\begin{document}
\maketitle
\noindent\textit{Editorial note.} All delivered mathematical text is the author's own, copied from the pinned Stacks Project chapter file etale-cohomology.tex; no statement or proof step was rewritten, reordered or removed.\par\medskip
\noindent\textit{Editorial note.} The counted proof is the author's complete argument: the proof is written out in full: the cofiltered limit presentation of the base change, the rewrite of the compared fibre product as a limit, and the application of the colimit theorem for cohomology.\par\medskip
\noindent\textit{Editorial note.} The label lemma-check-stalks-better names the counted statement, and the proof environment immediately following it in the source is delivered verbatim as its proof.\par\medskip
\noindent\textit{Editorial note.} The counted statement is carried by the author's own lemma environment in the chapter source.\par\medskip
\noindent\textit{Editorial note.} The statement's two conditions are displayed formulas of the author, delivered verbatim; the single cited input is kept as an external pointer with its canonical Stacks tag.\par\medskip
\noindent\textit{Editorial note.} The preamble is the chapter's own native preamble excerpt already used for this source: the Xy-pic setup (native preamble lines 14-18) and the native theorem declarations with the mathematical macros (native preamble lines 154-217).\par\medskip
\noindent\textit{Allowed standard prerequisites (editorial).}\begin{itemize}
\item Strict henselisations and the identification of the fibre products. Same chapter, the standing notation O\textasciicircum{}sh\_\{X,x\} and O\textasciicircum{}sh\_\{S,s\} and the fibre products appearing in both conditions. Supplies the two schemes whose cohomology groups are compared in the statement.
\item Cohomology of a cofiltered limit of schemes. Same chapter, the theorem computing the cohomology of the limit as the colimit of the cohomologies along a cofiltered limit with the appropriate hypotheses. Supplies the passage from the limit presentation of the fibre product to the conclusion.
\end{itemize}\medskip
\ifdefined\hypertarget\hypertarget{sourcepointer-1}{}\fi
\noindent\textit{Source reference.} Original reference 09YQ: sources/stacks-etale-cohomology/pinned-original-source.tex, source lines 7400--7400. Prerequisite pointer invoked at proof line(s) 16578; it names a standard fact of the same chapter that the counted argument uses, and every step of the counted proof built on it is retained verbatim in the delivered body.\par\medskip
\begin{lemma}
\label{lemma-check-stalks-better}
Let $f : X \to S$ be a morphism of schemes.
Let $\overline{x}$ be a geometric point of $X$ with image $\overline{s}$ in $S$.
Let $\Spec(K) \to \Spec(\mathcal{O}^{sh}_{S, \overline{s}})$
be a morphism with $K$ a separably closed field. Let $\mathcal{F}$ be an
abelian sheaf on $\Spec(K)_\etale$. Let $q \geq 0$. The following are
equivalent
\begin{enumerate}
\item
$H^q(\Spec(\mathcal{O}^{sh}_{X, \overline{x}}) \times_S \Spec(K), \mathcal{F}) =
H^q(\Spec(\mathcal{O}^{sh}_{S, \overline{s}}) \times_S \Spec(K), \mathcal{F})$
\item
$H^q(\Spec(\mathcal{O}^{sh}_{X, \overline{x}})
\times_{\Spec(\mathcal{O}^{sh}_{S, \overline{s}})} \Spec(K), \mathcal{F}) =
H^q(\Spec(K), \mathcal{F})$
\end{enumerate}
\end{lemma}

\begin{proof}
Observe that $\Spec(K) \times_S \Spec(\mathcal{O}^{sh}_{S, \overline{s}})$
is the spectrum of a filtered colimit of \'etale algebras over $K$.
Since $K$ is separably closed, each \'etale $K$-algebra
is a finite product of copies of $K$. Thus we can write
$$
\Spec(K) \times_S \Spec(\mathcal{O}^{sh}_{S, \overline{s}}) =
\lim_{i \in I} \coprod\nolimits_{a \in A_i} \Spec(K)
$$
as a cofiltered limit where each term is a disjoint union
of copies of $\Spec(K)$ over a finite set $A_i$.
Note that $A_i$ is nonempty as we are given
$\Spec(K) \to \Spec(\mathcal{O}^{sh}_{S, \overline{s}})$.
It follows that
\begin{align*}
\Spec(\mathcal{O}^{sh}_{X, \overline{x}}) \times_S \Spec(K)
& =
\Spec(\mathcal{O}^{sh}_{X, \overline{x}})
\times_{\Spec(\mathcal{O}^{sh}_{S, \overline{s}})}
\left(
\Spec(\mathcal{O}^{sh}_{S, \overline{s}})
\times_S \Spec(K)\right) \\
& =
\lim_{i \in I} \coprod\nolimits_{a \in A_i}
\Spec(\mathcal{O}^{sh}_{X, \overline{x}})
\times_{\Spec(\mathcal{O}^{sh}_{S, \overline{s}})} \Spec(K)
\end{align*}
Since taking cohomology in our setting commutes with limits
of schemes (Theorem \ref{theorem-colimit}) we conclude.
\end{proof}

\end{document}
